
\documentclass{article}
\usepackage{colortbl}
\usepackage{makecell}
\usepackage{multirow}
\usepackage{supertabular}

\begin{document}

\newcounter{utterance}

\centering \large Interaction Transcript for game `matchit\_ascii', experiment `similar\_grid\_2', episode 0 with nemotron{-}3{-}ultra.
\vspace{24pt}

{ \footnotesize  \setcounter{utterance}{1}
\setlength{\tabcolsep}{0pt}
\begin{supertabular}{c@{$\;$}|p{.15\linewidth}@{}p{.15\linewidth}p{.15\linewidth}p{.15\linewidth}p{.15\linewidth}p{.15\linewidth}}
    \# & $\;$A & \multicolumn{4}{c}{Game Master} & $\;\:$B\\
    \hline

    \theutterance \stepcounter{utterance}  
    & & \multicolumn{4}{p{0.6\linewidth}}{
        \cellcolor[rgb]{0.9,0.9,0.9}{
            \makecell[{{p{\linewidth}}}]{
                \texttt{\tiny{[P1$\langle$GM]}}
                \texttt{Vous participez à un jeu de devinettes collaboratif.} \\
\texttt{L'objectif est de déterminer si cette grille et une autre image que je suis le seul à voir sont identiques.} \\
\texttt{Veuillez d'abord décrire votre grille. Ensuite, je fournirai ma description et nous pourrons nous poser mutuellement des questions sur les grilles pour déterminer si elles sont identiques. Commencez maintenant votre brève description de la grille par "DESCRIPTION :" suivi de la description. N'ajoutez rien d'autre. Voici votre grille :} \\
\texttt{▢ X X X X} \\
\texttt{▢ X X X X} \\
\texttt{▢ X X X X} \\
\texttt{▢ X X X X} \\
\texttt{▢ ▢ ▢ ▢ ▢} \\
            }
        }
    }
    & \\ \\

    \theutterance \stepcounter{utterance}  
    & \multicolumn{4}{p{0.6\linewidth}}{
        \cellcolor[rgb]{0.9,0.9,0.9}{
            \makecell[{{p{\linewidth}}}]{
                \texttt{\tiny{[P1$\rangle$GM]}}
                \texttt{DESCRIPTION : Grille de 5 lignes et 5 colonnes. Les 4 premières lignes : première case vide (▢), quatre cases avec X. La 5ème ligne : cinq cases vides (▢).} \\
            }
        }
    }
    & & \\ \\

    \theutterance \stepcounter{utterance}  
    & & & \multicolumn{2}{p{0.3\linewidth}}{
        \cellcolor[rgb]{0.9,0.9,0.9}{
            \makecell[{{p{\linewidth}}}]{
                \texttt{\tiny{[GM$|$GM]}}
                \texttt{continue} \\
            }
        }
    }
    & & \\ \\

    \theutterance \stepcounter{utterance}  
    & & \multicolumn{4}{p{0.6\linewidth}}{
        \cellcolor[rgb]{0.9,0.9,0.9}{
            \makecell[{{p{\linewidth}}}]{
                \texttt{\tiny{[GM$\rangle$P2]}}
                \texttt{Vous participez à un jeu de devinettes collaboratif.} \\
\texttt{L'objectif est de déterminer si cette grille et une autre image que je suis le seul à voir sont identiques.} \\
\texttt{Veuillez d'abord décrire votre grille. Ensuite, je fournirai ma description et nous pourrons nous poser mutuellement des questions sur les grilles pour déterminer si elles sont identiques. Commencez maintenant votre brève description de la grille par "DESCRIPTION :" suivi de la description. N'ajoutez rien d'autre. Voici votre grille :} \\
\texttt{X X X X X} \\
\texttt{▢ X X X ▢} \\
\texttt{▢ X X X X} \\
\texttt{▢ X X X X} \\
\texttt{▢ ▢ ▢ ▢ ▢} \\
            }
        }
    }
    & \\ \\

    \theutterance \stepcounter{utterance}  
    & & & \multicolumn{4}{p{0.6\linewidth}}{
        \cellcolor[rgb]{0.9,0.9,0.9}{
            \makecell[{{p{\linewidth}}}]{
                \texttt{\tiny{[GM$\langle$P2]}}
                \texttt{DESCRIPTION : Grille 5x5. Première ligne : cinq X. Deuxième ligne : case vide, trois X, case vide. Troisième et quatrième lignes : case vide suivies de quatre X. Cinquième ligne : cinq cases vides.} \\
            }
        }
    }
     \\ \\

    \theutterance \stepcounter{utterance}  
    & & & \multicolumn{2}{p{0.3\linewidth}}{
        \cellcolor[rgb]{0.9,0.9,0.9}{
            \makecell[{{p{\linewidth}}}]{
                \texttt{\tiny{[GM$|$GM]}}
                \texttt{continue} \\
            }
        }
    }
    & & \\ \\

    \theutterance \stepcounter{utterance}  
    & & \multicolumn{4}{p{0.6\linewidth}}{
        \cellcolor[rgb]{0.9,0.9,0.9}{
            \makecell[{{p{\linewidth}}}]{
                \texttt{\tiny{[GM$\rangle$P2]}}
                \texttt{Voici ma DESCRIPTIONDESCRIPTION : Grille de 5 lignes et 5 colonnes. Les 4 premières lignes : première case vide (▢), quatre cases avec X. La 5ème ligne : cinq cases vides (▢).} \\
\texttt{Posez maintenant une question afin de découvrir de nouveaux aspects de ma grille qui pourraient être différents de votre grille. Commencez par "QUESTION :" et n'ajoutez rien d'autre.} \\
            }
        }
    }
    & \\ \\

    \theutterance \stepcounter{utterance}  
    & & & \multicolumn{4}{p{0.6\linewidth}}{
        \cellcolor[rgb]{0.9,0.9,0.9}{
            \makecell[{{p{\linewidth}}}]{
                \texttt{\tiny{[GM$\langle$P2]}}
                \texttt{QUESTION : Sur votre première ligne, la première case est{-}elle vide ou contient{-}elle un X ?} \\
            }
        }
    }
     \\ \\

    \theutterance \stepcounter{utterance}  
    & & & \multicolumn{2}{p{0.3\linewidth}}{
        \cellcolor[rgb]{0.9,0.9,0.9}{
            \makecell[{{p{\linewidth}}}]{
                \texttt{\tiny{[GM$|$GM]}}
                \texttt{continue} \\
            }
        }
    }
    & & \\ \\

    \theutterance \stepcounter{utterance}  
    & & \multicolumn{4}{p{0.6\linewidth}}{
        \cellcolor[rgb]{0.9,0.9,0.9}{
            \makecell[{{p{\linewidth}}}]{
                \texttt{\tiny{[P1$\langle$GM]}}
                \texttt{Voici ma DESCRIPTIONDESCRIPTION : Grille 5x5. Première ligne : cinq X. Deuxième ligne : case vide, trois X, case vide. Troisième et quatrième lignes : case vide suivies de quatre X. Cinquième ligne : cinq cases vides.} \\
\texttt{QUESTION : Sur votre première ligne, la première case est{-}elle vide ou contient{-}elle un X ?Commencez votre réponse par "RÉPONSE :" et n'ajoutez rien d'autre.} \\
            }
        }
    }
    & \\ \\

    \theutterance \stepcounter{utterance}  
    & \multicolumn{4}{p{0.6\linewidth}}{
        \cellcolor[rgb]{0.9,0.9,0.9}{
            \makecell[{{p{\linewidth}}}]{
                \texttt{\tiny{[P1$\rangle$GM]}}
                \texttt{RÉPONSE : La première case de la première ligne est vide (▢).} \\
            }
        }
    }
    & & \\ \\

    \theutterance \stepcounter{utterance}  
    & & & \multicolumn{2}{p{0.3\linewidth}}{
        \cellcolor[rgb]{0.9,0.9,0.9}{
            \makecell[{{p{\linewidth}}}]{
                \texttt{\tiny{[GM$|$GM]}}
                \texttt{continue} \\
            }
        }
    }
    & & \\ \\

    \theutterance \stepcounter{utterance}  
    & & \multicolumn{4}{p{0.6\linewidth}}{
        \cellcolor[rgb]{0.9,0.9,0.9}{
            \makecell[{{p{\linewidth}}}]{
                \texttt{\tiny{[P1$\langle$GM]}}
                \texttt{Posez maintenant une question afin de découvrir de nouveaux aspects de ma grille qui pourraient être différents de votre grille. Commencez par "QUESTION :" et n'ajoutez rien d'autre.} \\
            }
        }
    }
    & \\ \\

    \theutterance \stepcounter{utterance}  
    & \multicolumn{4}{p{0.6\linewidth}}{
        \cellcolor[rgb]{0.9,0.9,0.9}{
            \makecell[{{p{\linewidth}}}]{
                \texttt{\tiny{[P1$\rangle$GM]}}
                \texttt{QUESTION : Sur votre deuxième ligne, les cases 2, 3 et 4 sont{-}elles toutes des X ?} \\
            }
        }
    }
    & & \\ \\

    \theutterance \stepcounter{utterance}  
    & & & \multicolumn{2}{p{0.3\linewidth}}{
        \cellcolor[rgb]{0.9,0.9,0.9}{
            \makecell[{{p{\linewidth}}}]{
                \texttt{\tiny{[GM$|$GM]}}
                \texttt{continue} \\
            }
        }
    }
    & & \\ \\

    \theutterance \stepcounter{utterance}  
    & & \multicolumn{4}{p{0.6\linewidth}}{
        \cellcolor[rgb]{0.9,0.9,0.9}{
            \makecell[{{p{\linewidth}}}]{
                \texttt{\tiny{[GM$\rangle$P2]}}
                \texttt{RÉPONSE : La première case de la première ligne est vide (▢).} \\
\texttt{QUESTION : Sur votre deuxième ligne, les cases 2, 3 et 4 sont{-}elles toutes des X ?Commencez votre réponse par "RÉPONSE :" et n'ajoutez rien d'autre.} \\
            }
        }
    }
    & \\ \\

    \theutterance \stepcounter{utterance}  
    & & & \multicolumn{4}{p{0.6\linewidth}}{
        \cellcolor[rgb]{0.9,0.9,0.9}{
            \makecell[{{p{\linewidth}}}]{
                \texttt{\tiny{[GM$\langle$P2]}}
                \texttt{RÉPONSE : Oui, les cases 2, 3 et 4 de ma deuxième ligne sont toutes des X.} \\
\texttt{QUESTION : Sur votre deuxième ligne, la cinquième case (colonne 5) est{-}elle un X ou une case vide ?} \\
            }
        }
    }
     \\ \\

    \theutterance \stepcounter{utterance}  
    & & & \multicolumn{2}{p{0.3\linewidth}}{
        \cellcolor[rgb]{0.9,0.9,0.9}{
            \makecell[{{p{\linewidth}}}]{
                \texttt{\tiny{[GM$|$GM]}}
                \texttt{continue} \\
            }
        }
    }
    & & \\ \\

    \theutterance \stepcounter{utterance}  
    & & \multicolumn{4}{p{0.6\linewidth}}{
        \cellcolor[rgb]{0.9,0.9,0.9}{
            \makecell[{{p{\linewidth}}}]{
                \texttt{\tiny{[GM$\rangle$P2]}}
                \texttt{Posez maintenant une question afin de découvrir de nouveaux aspects de ma grille qui pourraient être différents de votre grille. Commencez par "QUESTION :" et n'ajoutez rien d'autre.} \\
            }
        }
    }
    & \\ \\

    \theutterance \stepcounter{utterance}  
    & & & \multicolumn{4}{p{0.6\linewidth}}{
        \cellcolor[rgb]{0.9,0.9,0.9}{
            \makecell[{{p{\linewidth}}}]{
                \texttt{\tiny{[GM$\langle$P2]}}
                \texttt{QUESTION : Sur votre troisième ligne, la première case est{-}elle vide et les quatre suivantes sont{-}elles des X ?} \\
            }
        }
    }
     \\ \\

    \theutterance \stepcounter{utterance}  
    & & & \multicolumn{2}{p{0.3\linewidth}}{
        \cellcolor[rgb]{0.9,0.9,0.9}{
            \makecell[{{p{\linewidth}}}]{
                \texttt{\tiny{[GM$|$GM]}}
                \texttt{continue} \\
            }
        }
    }
    & & \\ \\

    \theutterance \stepcounter{utterance}  
    & & \multicolumn{4}{p{0.6\linewidth}}{
        \cellcolor[rgb]{0.9,0.9,0.9}{
            \makecell[{{p{\linewidth}}}]{
                \texttt{\tiny{[P1$\langle$GM]}}
                \texttt{RÉPONSE : Oui, les cases 2, 3 et 4 de ma deuxième ligne sont toutes des X.} \\
\texttt{QUESTION : Sur votre deuxième ligne, la cinquième case (colonne 5) est{-}elle un X ou une case vide ?} \\
\texttt{QUESTION : Sur votre troisième ligne, la première case est{-}elle vide et les quatre suivantes sont{-}elles des X ?Commencez votre réponse par "RÉPONSE :" et n'ajoutez rien d'autre.} \\
            }
        }
    }
    & \\ \\

    \theutterance \stepcounter{utterance}  
    & \multicolumn{4}{p{0.6\linewidth}}{
        \cellcolor[rgb]{0.9,0.9,0.9}{
            \makecell[{{p{\linewidth}}}]{
                \texttt{\tiny{[P1$\rangle$GM]}}
                \texttt{RÉPONSE : Sur ma deuxième ligne, la cinquième case (colonne 5) est un X. Sur ma troisième ligne, la première case est vide (▢) et les quatre suivantes sont des X.} \\
            }
        }
    }
    & & \\ \\

    \theutterance \stepcounter{utterance}  
    & & & \multicolumn{2}{p{0.3\linewidth}}{
        \cellcolor[rgb]{0.9,0.9,0.9}{
            \makecell[{{p{\linewidth}}}]{
                \texttt{\tiny{[GM$|$GM]}}
                \texttt{continue} \\
            }
        }
    }
    & & \\ \\

    \theutterance \stepcounter{utterance}  
    & & \multicolumn{4}{p{0.6\linewidth}}{
        \cellcolor[rgb]{0.9,0.9,0.9}{
            \makecell[{{p{\linewidth}}}]{
                \texttt{\tiny{[P1$\langle$GM]}}
                \texttt{Posez maintenant une question afin de découvrir de nouveaux aspects de ma grille qui pourraient être différents de votre grille. Commencez par "QUESTION :" et n'ajoutez rien d'autre.} \\
            }
        }
    }
    & \\ \\

    \theutterance \stepcounter{utterance}  
    & \multicolumn{4}{p{0.6\linewidth}}{
        \cellcolor[rgb]{0.9,0.9,0.9}{
            \makecell[{{p{\linewidth}}}]{
                \texttt{\tiny{[P1$\rangle$GM]}}
                \texttt{QUESTION : Sur votre première ligne, les cinq cases sont{-}elles toutes des X ?} \\
            }
        }
    }
    & & \\ \\

    \theutterance \stepcounter{utterance}  
    & & & \multicolumn{2}{p{0.3\linewidth}}{
        \cellcolor[rgb]{0.9,0.9,0.9}{
            \makecell[{{p{\linewidth}}}]{
                \texttt{\tiny{[GM$|$GM]}}
                \texttt{continue} \\
            }
        }
    }
    & & \\ \\

    \theutterance \stepcounter{utterance}  
    & & \multicolumn{4}{p{0.6\linewidth}}{
        \cellcolor[rgb]{0.9,0.9,0.9}{
            \makecell[{{p{\linewidth}}}]{
                \texttt{\tiny{[GM$\rangle$P2]}}
                \texttt{RÉPONSE : Sur ma deuxième ligne, la cinquième case (colonne 5) est un X. Sur ma troisième ligne, la première case est vide (▢) et les quatre suivantes sont des X.} \\
\texttt{QUESTION : Sur votre première ligne, les cinq cases sont{-}elles toutes des X ?Commencez votre réponse par "RÉPONSE :" et n'ajoutez rien d'autre.} \\
            }
        }
    }
    & \\ \\

    \theutterance \stepcounter{utterance}  
    & & & \multicolumn{4}{p{0.6\linewidth}}{
        \cellcolor[rgb]{0.9,0.9,0.9}{
            \makecell[{{p{\linewidth}}}]{
                \texttt{\tiny{[GM$\langle$P2]}}
                \texttt{RÉPONSE : Oui, sur ma première ligne, les cinq cases sont toutes des X.} \\
\texttt{QUESTION : Sur votre quatrième ligne, la première case est{-}elle vide et les quatre suivantes sont{-}elles des X ?} \\
            }
        }
    }
     \\ \\

    \theutterance \stepcounter{utterance}  
    & & & \multicolumn{2}{p{0.3\linewidth}}{
        \cellcolor[rgb]{0.9,0.9,0.9}{
            \makecell[{{p{\linewidth}}}]{
                \texttt{\tiny{[GM$|$GM]}}
                \texttt{continue} \\
            }
        }
    }
    & & \\ \\

    \theutterance \stepcounter{utterance}  
    & & \multicolumn{4}{p{0.6\linewidth}}{
        \cellcolor[rgb]{0.9,0.9,0.9}{
            \makecell[{{p{\linewidth}}}]{
                \texttt{\tiny{[GM$\rangle$P2]}}
                \texttt{Posez maintenant une question afin de découvrir de nouveaux aspects de ma grille qui pourraient être différents de votre grille. Commencez par "QUESTION :" et n'ajoutez rien d'autre.} \\
            }
        }
    }
    & \\ \\

    \theutterance \stepcounter{utterance}  
    & & & \multicolumn{4}{p{0.6\linewidth}}{
        \cellcolor[rgb]{0.9,0.9,0.9}{
            \makecell[{{p{\linewidth}}}]{
                \texttt{\tiny{[GM$\langle$P2]}}
                \texttt{QUESTION : Sur votre quatrième ligne, la première case est{-}elle vide (▢) et les quatre cases suivantes sont{-}elles des X ?} \\
            }
        }
    }
     \\ \\

    \theutterance \stepcounter{utterance}  
    & & & \multicolumn{2}{p{0.3\linewidth}}{
        \cellcolor[rgb]{0.9,0.9,0.9}{
            \makecell[{{p{\linewidth}}}]{
                \texttt{\tiny{[GM$|$GM]}}
                \texttt{continue} \\
            }
        }
    }
    & & \\ \\

    \theutterance \stepcounter{utterance}  
    & & \multicolumn{4}{p{0.6\linewidth}}{
        \cellcolor[rgb]{0.9,0.9,0.9}{
            \makecell[{{p{\linewidth}}}]{
                \texttt{\tiny{[P1$\langle$GM]}}
                \texttt{RÉPONSE : Oui, sur ma première ligne, les cinq cases sont toutes des X.} \\
\texttt{QUESTION : Sur votre quatrième ligne, la première case est{-}elle vide et les quatre suivantes sont{-}elles des X ?} \\
\texttt{QUESTION : Sur votre quatrième ligne, la première case est{-}elle vide (▢) et les quatre cases suivantes sont{-}elles des X ?Commencez votre réponse par "RÉPONSE :" et n'ajoutez rien d'autre.} \\
            }
        }
    }
    & \\ \\

    \theutterance \stepcounter{utterance}  
    & \multicolumn{4}{p{0.6\linewidth}}{
        \cellcolor[rgb]{0.9,0.9,0.9}{
            \makecell[{{p{\linewidth}}}]{
                \texttt{\tiny{[P1$\rangle$GM]}}
                \texttt{RÉPONSE : Oui, sur ma quatrième ligne, la première case est vide (▢) et les quatre cases suivantes sont des X.} \\
            }
        }
    }
    & & \\ \\

    \theutterance \stepcounter{utterance}  
    & & & \multicolumn{2}{p{0.3\linewidth}}{
        \cellcolor[rgb]{0.9,0.9,0.9}{
            \makecell[{{p{\linewidth}}}]{
                \texttt{\tiny{[GM$|$GM]}}
                \texttt{continue} \\
            }
        }
    }
    & & \\ \\

    \theutterance \stepcounter{utterance}  
    & & \multicolumn{4}{p{0.6\linewidth}}{
        \cellcolor[rgb]{0.9,0.9,0.9}{
            \makecell[{{p{\linewidth}}}]{
                \texttt{\tiny{[P1$\langle$GM]}}
                \texttt{Posez maintenant une question afin de découvrir de nouveaux aspects de ma grille qui pourraient être différents de votre grille. Commencez par "QUESTION :" et n'ajoutez rien d'autre.} \\
            }
        }
    }
    & \\ \\

    \theutterance \stepcounter{utterance}  
    & \multicolumn{4}{p{0.6\linewidth}}{
        \cellcolor[rgb]{0.9,0.9,0.9}{
            \makecell[{{p{\linewidth}}}]{
                \texttt{\tiny{[P1$\rangle$GM]}}
                \texttt{QUESTION : Sur votre deuxième ligne, la première case est{-}elle vide (▢) et la cinquième case est{-}elle vide (▢) ?} \\
            }
        }
    }
    & & \\ \\

    \theutterance \stepcounter{utterance}  
    & & & \multicolumn{2}{p{0.3\linewidth}}{
        \cellcolor[rgb]{0.9,0.9,0.9}{
            \makecell[{{p{\linewidth}}}]{
                \texttt{\tiny{[GM$|$GM]}}
                \texttt{continue} \\
            }
        }
    }
    & & \\ \\

    \theutterance \stepcounter{utterance}  
    & & \multicolumn{4}{p{0.6\linewidth}}{
        \cellcolor[rgb]{0.9,0.9,0.9}{
            \makecell[{{p{\linewidth}}}]{
                \texttt{\tiny{[GM$\rangle$P2]}}
                \texttt{RÉPONSE : Oui, sur ma quatrième ligne, la première case est vide (▢) et les quatre cases suivantes sont des X.} \\
\texttt{QUESTION : Sur votre deuxième ligne, la première case est{-}elle vide (▢) et la cinquième case est{-}elle vide (▢) ?Commencez votre réponse par "RÉPONSE :" et n'ajoutez rien d'autre.} \\
            }
        }
    }
    & \\ \\

    \theutterance \stepcounter{utterance}  
    & & & \multicolumn{4}{p{0.6\linewidth}}{
        \cellcolor[rgb]{0.9,0.9,0.9}{
            \makecell[{{p{\linewidth}}}]{
                \texttt{\tiny{[GM$\langle$P2]}}
                \texttt{RÉPONSE : Oui, sur ma deuxième ligne, la première case est vide (▢) et la cinquième case est vide (▢).} \\
\texttt{QUESTION : Sur votre cinquième ligne, les cinq cases sont{-}elles toutes vides (▢) ?} \\
            }
        }
    }
     \\ \\

    \theutterance \stepcounter{utterance}  
    & & & \multicolumn{2}{p{0.3\linewidth}}{
        \cellcolor[rgb]{0.9,0.9,0.9}{
            \makecell[{{p{\linewidth}}}]{
                \texttt{\tiny{[GM$|$GM]}}
                \texttt{continue} \\
            }
        }
    }
    & & \\ \\

    \theutterance \stepcounter{utterance}  
    & & \multicolumn{4}{p{0.6\linewidth}}{
        \cellcolor[rgb]{0.9,0.9,0.9}{
            \makecell[{{p{\linewidth}}}]{
                \texttt{\tiny{[GM$\rangle$P2]}}
                \texttt{Prenez maintenant une décision. Qu’en pensez{-}vous : votre grille et l’autre grille décrite sont{-}elles identiques ? Écrivez "DÉCISION : même grille" si vous pensez qu’il s’agit de la même grille ou "DÉCISION : grilles différentes" si vous pensez qu’il s’agit de grilles différentes. N’ajoutez rien d’autre.} \\
            }
        }
    }
    & \\ \\

    \theutterance \stepcounter{utterance}  
    & & & \multicolumn{4}{p{0.6\linewidth}}{
        \cellcolor[rgb]{0.9,0.9,0.9}{
            \makecell[{{p{\linewidth}}}]{
                \texttt{\tiny{[GM$\langle$P2]}}
                \texttt{DÉCISION : grilles différentes} \\
            }
        }
    }
     \\ \\

    \theutterance \stepcounter{utterance}  
    & & & \multicolumn{2}{p{0.3\linewidth}}{
        \cellcolor[rgb]{0.9,0.9,0.9}{
            \makecell[{{p{\linewidth}}}]{
                \texttt{\tiny{[GM$|$GM]}}
                \texttt{continue} \\
            }
        }
    }
    & & \\ \\

    \theutterance \stepcounter{utterance}  
    & & & \multicolumn{2}{p{0.3\linewidth}}{
        \cellcolor[rgb]{0.9,0.9,0.9}{
            \makecell[{{p{\linewidth}}}]{
                \texttt{\tiny{[GM$|$GM]}}
                \texttt{success} \\
            }
        }
    }
    & & \\ \\

    \theutterance \stepcounter{utterance}  
    & & \multicolumn{4}{p{0.6\linewidth}}{
        \cellcolor[rgb]{0.9,0.9,0.9}{
            \makecell[{{p{\linewidth}}}]{
                \texttt{\tiny{[P1$\langle$GM]}}
                \texttt{RÉPONSE : Oui, sur ma deuxième ligne, la première case est vide (▢) et la cinquième case est vide (▢).} \\
\texttt{QUESTION : Sur votre cinquième ligne, les cinq cases sont{-}elles toutes vides (▢) ?} \\
\texttt{Prenez maintenant une décision. Qu’en pensez{-}vous : votre grille et l’autre grille décrite sont{-}elles identiques ? Écrivez "DÉCISION : même grille" si vous pensez qu’il s’agit de la même grille ou "DÉCISION : grilles différentes" si vous pensez qu’il s’agit de grilles différentes. N’ajoutez rien d’autre.} \\
            }
        }
    }
    & \\ \\

    \theutterance \stepcounter{utterance}  
    & \multicolumn{4}{p{0.6\linewidth}}{
        \cellcolor[rgb]{0.9,0.9,0.9}{
            \makecell[{{p{\linewidth}}}]{
                \texttt{\tiny{[P1$\rangle$GM]}}
                \texttt{DÉCISION : grilles différentes} \\
            }
        }
    }
    & & \\ \\

    \theutterance \stepcounter{utterance}  
    & & & \multicolumn{2}{p{0.3\linewidth}}{
        \cellcolor[rgb]{0.9,0.9,0.9}{
            \makecell[{{p{\linewidth}}}]{
                \texttt{\tiny{[GM$|$GM]}}
                \texttt{continue} \\
            }
        }
    }
    & & \\ \\

    \theutterance \stepcounter{utterance}  
    & & & \multicolumn{2}{p{0.3\linewidth}}{
        \cellcolor[rgb]{0.9,0.9,0.9}{
            \makecell[{{p{\linewidth}}}]{
                \texttt{\tiny{[GM$|$GM]}}
                \texttt{success} \\
            }
        }
    }
    & & \\ \\

\end{supertabular}
}

\end{document}
